%%%%%%%%%%%%%%%%%%%%%%%%%%%%%%%%%%%%%%%%%%%%%%%%%%%%%%%%%%%%%%%%%%%
\section{Background}
\label{sec:background}

%---------------------------------------------------------------
\subsection{Linear Programming Duality and Sensitivity}
\label{sec:lp-duality}
\begin{itemize}
	\item State the primal LP $\min c^\top x$ s.t. $Ax \ge b,\ x \ge 0$ and its dual.
	\item Introduce the basis: with $m$ equality constraints, a basic solution
	has exactly $m$ basic variables and the basis matrix is invertible.
\end{itemize}

\paragraph{Basic vs.\ non-basic variables and complementary slackness.}
\begin{itemize}
	\item A basic variable has \ac{RC} zero by definition; a variable with
	non-zero \ac{RC} is provably non-basic.
	\item Complementary slackness links primal activity and dual constraints.
\end{itemize}

\paragraph{Shadow price.}
\begin{itemize}
	\item The \ac{SP} of a constraint is the marginal change of the objective
	per unit relaxation of that constraint's right-hand side.
	\item Valid only within the basis-stable interval.
\end{itemize}

\paragraph{Reduced cost.}
\begin{itemize}
	\item The \ac{RC} of a (currently non-basic) variable is the amount by which
	its objective coefficient must improve before it enters the optimal basis at
	a positive level.
	\item Central reading: ``how much cheaper must this technology become to be
	built''.
\end{itemize}

%---------------------------------------------------------------
\subsection{Solving Linear Programs}
\label{sec:solving-lps}

\paragraph{Simplex (primal and dual).}
\begin{itemize}
	\item Vertex-following methods.
	\item Always terminate at a basic (vertex) solution, hence return a basis
	and well-defined \ac{RC}/\ac{SP}.
\end{itemize}

\paragraph{Barrier / interior point.}
\begin{itemize}
	\item Follow the central path to an interior optimum.
	\item On their own they return a non-vertex solution, so basis-dependent
	duals are not directly available.
\end{itemize}

\paragraph{Crossover.}
\begin{itemize}
	\item Pushes a barrier solution to an adjacent vertex so a basis -- and
	therefore exact \ac{RC}/\ac{SP} -- becomes available.
	\item Required for the VBasis flag used in Section~\ref{sec:methodology}.
\end{itemize}

\paragraph{Presolve and scaling.}
\begin{itemize}
	\item Presolve aggregates/substitutes constraints and variables for speed,
	which can make certain constraint duals unrecoverable in the original space.
	\item Scaling rescales rows/columns and must be undone to report duals in
	input units.
	\item Both effects motivate the robustness study in Section~\ref{sec:results}.
\end{itemize}

%---------------------------------------------------------------
\subsection{Boundary Information in Established Energy System Models}
\label{sec:related-work}
tbd
